\documentclass[11pt]{article}
\usepackage[a4paper,margin=1in]{geometry}
\usepackage{amsmath,amssymb,amsthm,booktabs,hyperref,xcolor,bm}
\newcommand{\R}{\mathcal R}
\newcommand{\norm}[1]{\lVert #1 \rVert}
\title{Testing the Excess-Risk Certificate Axis by Axis\\
       on the Rectified-Flow Family}
\author{thermau5}
\date{\today}
\begin{document}
\maketitle

\begin{abstract}
The generic certificate bounds excess terminal risk as
$Q = Q_0 + \tfrac{c_{\mathrm{stat}}}{n}\!\int\! \tfrac{a}{\rho}
      + c_{\mathrm{disc}}\!\int\! \tfrac{d}{m^{p}}$,
with optimal sampling density $m^\star\!\propto d^{1/(p+1)}$. Each design
choice enters through a specific term (Table~\ref{tab:axes}). Prior locked
results varied the schedule $m$ and solver $s$ at a fixed EDM path; here we
test every axis we can on a \emph{new} model family, Rectified Flow (RF),
to see what generalizes. The two headline findings: (i) on the
\textbf{path} axis, the reflow ladder 1/2/3-RF produces the
floor-vs-defect \emph{crossover cascade} the certificate predicts --- no
single path wins at all budgets; (ii) on the \textbf{schedule} axis, a
\emph{parameter-free} certificate-calibrated grid wins or ties the
\emph{native default} of every configuration we test --- four EDM solver
cores and four path families (EDM, RF, VP-SDE, OT-CFM) --- the only losses
being at the lowest budgets (EMS at NFE 5--8, whose schedule is co-optimized
with the solver; OT-CFM at NFE 5; Table~\ref{tab:dominance}). Where we can
search exhaustively, the closed form is statistically indistinguishable from
the \emph{FID-optimal} schedule on 1-RF, with $0.7$--$1.3$ FID of quantified
headroom on EDM/Heun at NFE 9--13; against \emph{optimized} schedule
baselines (per-NFE tuned $\rho$, the GITS dynamic program, a surrogate-DP;
Table~\ref{tab:schedaxis}) it is the only scheduler that is never
catastrophic in any group, at $10$--$500\times$ less preparation compute.
It also transfers across datasets --- on FFHQ-64 it beats the Karras default
at every budget with the CIFAR-validated exponents unchanged. All FID is
Clean-FID 10k, 3 seeds (FFHQ: EDM's FID implementation, both arms).
\end{abstract}

\section{The hierarchy}
\label{sec:hierarchy}
The certificate decomposes by design axis: each knob controls one part of
the bound. This is the organizing structure of the whole study.

\begin{table}[h]\centering
\caption{The certificate's axes. Each row is one experiment in this report.}
\label{tab:axes}
\begin{tabular}{llll}
\toprule
Axis & Knob & Enters via & Tested in \\
\midrule
Solver   & $s$ & local defect $d$            & \S\ref{sec:solver} (Table~\ref{tab:edm})\\
Risk     & $\R$& sensitivity in $a,d$        & \S\ref{sec:risk} (Table~\ref{tab:risk})\\
Path     & $P$ & floor $Q_0$ \emph{and} $d$  & \S\ref{sec:path} (Table~\ref{tab:path})\\
Schedule & $m$ & disc.\ term $\int d/m^{p}$  & \S\ref{sec:sched} (Table~\ref{tab:dominance})\\
\bottomrule
\end{tabular}
\end{table}

\noindent\textbf{Results at a glance.}
\begin{itemize}\itemsep2pt
\item[$s$:] DPM-Solver-v3 (EMS) is the strongest solver core --- best floor
($\approx4.35$), winning at mid-to-high NFE on the fixed EDM path --- a
Level-1 advance.
\item[$\R$:] The directly-measurable feature-space $W_2$ risk tracks FID
in rank $5/5$ --- it is the FID-faithful risk.
\item[$P$:] Reflow trades the floor $Q_0$ for the defect $d$; the optimal
reflow count is \emph{NFE-dependent} (3-RF at NFE 5, 2-RF at 8--64, 1-RF at
$\ge128$). Pre-registered crossover prediction confirmed.
\item[$m$:] The certificate's calibrated node placement
$m^\star\!\propto d^{1/(p+1)}$ ($p\!=\!1$, no metric feedback) wins or ties
the native default schedule of every solver core and path family tested,
except at the lowest budgets (EMS NFE 5--8, OT-CFM NFE 5;
Table~\ref{tab:dominance}); a free-node search cannot beat it on 1-RF
(FID-optimal there) and quantifies $0.7$--$1.3$ FID of headroom on EDM/Heun
at NFE 9--13; among optimized schedulers (Table~\ref{tab:schedaxis}) it is
the only one never catastrophic in any (path, solver) group; and it
transfers to FFHQ-64 with no re-tuning, winning every budget.
\end{itemize}

\section{Solver axis $s$: fixed-EDM frontier}
\label{sec:solver}
The bar the RF results are measured against. EDM CIFAR-10, locked protocol.

\begin{table}[h]\centering
\caption{Solver axis on a fixed EDM path. Clean-FID, lower better; bold best.}
\label{tab:edm}
\begin{tabular}{lcccccc}
\toprule
NFE & 5 & 8 & 12 & 18 & 32 & 64\\
\midrule
Ours (UniPC schedule) & 21.46 & 9.18 & 5.58 & 4.66 & 4.45 & 4.41\\
dpm\_solver\_v3 (EMS)  & \textbf{17.07} & \textbf{6.30} & \textbf{4.71} & \textbf{4.46} & \textbf{4.35} & \textbf{4.35}\\
\bottomrule
\end{tabular}
\end{table}

\noindent EMS beats our UniPC-schedule arm at every NFE here and has the best
floor ($\approx4.35$) on offer: a \emph{solver} ($s$) advance, not a schedule
advance, consistent with $s$ entering $d$. (It is not uniformly ahead of
\emph{every} schedule, though: at the two lowest budgets our calibrated Heun
schedule beats it --- Table~\ref{tab:dominance}, Heun vs EMS rows.)

\section{Risk axis $\R$: the FID-faithful risk}
\label{sec:risk}
Is the bounded risk the one we care about? We compare the directly-measured
$R_{\mathrm{disc}}=$ feature-$W_2$(K-step, 128-step ref) against FID on a
5-grid panel (3 seeds).

\begin{table}[h]\centering
\caption{$R_{\mathrm{disc}}$ vs FID. Identical rank ordering ($5/5$).}
\label{tab:risk}
\begin{tabular}{lcc}
\toprule
grid & $R_{\mathrm{disc}}$ & FID\\
\midrule
good    & 13.76  & 13.16\\
refined & 28.91  & 26.71\\
karras  & 50.43  & 44.90\\
forced  & 96.73  & 87.85\\
cluster & 282.87 & 270.30\\
\bottomrule
\end{tabular}
\end{table}

\noindent Feature-space $W_2$ \emph{is} the FID-faithful risk (the risk-axis
question is settled); pixel-$L_2$ scrambles this ranking.

\section{Path axis $P$: the reflow ladder}
\label{sec:path}
Hold solver and schedule fixed (Euler, uniform); vary only path straightness
via reflow 1/2/3-RF (same backbone, same $v$-parameterization). \emph{Pre-registered
prediction:} reflow lowers $d$ (helps low NFE) but raises $Q_0$ (hurts high
NFE), so the order must \emph{reverse} with NFE --- no single winner.

\begin{table}[h]\centering
\caption{Path ladder, uniform Euler. Clean-FID; bold best at each NFE.}
\label{tab:path}
\begin{tabular}{lcccccccc}
\toprule
NFE & 5 & 8 & 12 & 18 & 32 & 64 & 128 & 256\\
\midrule
1-RF & 38.59 & 20.21 & 13.76 & 10.57 & 8.04 & 6.53 & \textbf{5.85} & \textbf{5.56}\\
2-RF & 7.55 & \textbf{6.99} & \textbf{6.70} & \textbf{6.53} & \textbf{6.38} & \textbf{6.31} & 6.27 & 6.26\\
3-RF & \textbf{7.45} & 7.20 & 7.06 & 6.97 & 6.89 & 6.85 & 6.82 & 6.81\\
\bottomrule
\end{tabular}
\end{table}

\noindent\textbf{Confirmed.} The best path changes with budget ---
\,3-RF\,(NFE 5)\,$\to$\,2-RF\,(8--64)\,$\to$\,1-RF\,($\ge$128) --- a cascade
of three crossovers ($3\!\leftrightarrow\!2$ at ${\approx}6$,
$1\!\leftrightarrow\!3$ in $(32,64)$, $1\!\leftrightarrow\!2$ in $(64,128)$).
At NFE${\ge}128$ the order is the fully-reversed $1\!<\!2\!<\!3$-RF, matching
the published converged floors ($2.58\!<\!3.36\!<\!3.96$). This is exactly the
$Q_0$-vs-$d$ trade. \emph{Cross-method:} the straightened path beats the best
EDM solver (Table~\ref{tab:edm}) only at NFE 5 ($\approx7.5$ vs $17.07$); the
low-floor EDM model wins every larger budget.

\section{Schedule axis $m$: the dominance table}
\label{sec:sched}
The certificate prescribes a \emph{single} schedule,
$m^\star\!\propto d^{1/(p+1)}$ with $p\!=\!1$ and $d=\norm{\ddot x}$
calibrated on a held-out reference trajectory (no FID feedback, no tuned
scalar; \S\ref{app:protocol}). The question this section settles: does that
one prescription beat the \emph{native} default schedule of \emph{every}
solver core and \emph{every} path family --- is it optimal regardless of the
other axes? Table~\ref{tab:dominance} answers yes, with one documented
boundary.

\begin{table}[t]\centering
\caption{\textbf{Schedule dominance.} Each cell is paired Clean-FID,
\emph{our} $m^\star$ / that configuration's \emph{own default} schedule, at
matched NFE (3 seeds). \textbf{Bold} marks the better arm when it leads
beyond $2\sigma$ of the per-seed difference, or \emph{both} arms when they
tie within $2\sigma$ (joint-best). Across four EDM
solver cores and four path families, the \emph{only} losses are at the
lowest budgets --- EMS at NFE 5,\,8 and OT-CFM at NFE 5; everywhere else
(including EMS at NFE$\,\ge\,$12) our schedule ties or wins.}
\label{tab:dominance}
\footnotesize
\setlength{\tabcolsep}{4.5pt}
\begin{tabular}{l cccccc}
\toprule
$m^\star$\,/\,default & NFE\,5 & 8 & 12 & 18 & 32 & 64\\
\midrule
\multicolumn{7}{l}{\emph{Solver core $s$ --- fixed EDM path, vs.\ the core's own default schedule}}\\
Heun$^\dagger$ & \textbf{46.22}/338 & \textbf{13.41}/45.07 & \textbf{7.70}/11.64 & \textbf{5.33}/5.56 & \textbf{4.52}/\textbf{4.45} & \textbf{4.41}/\textbf{4.40}\\
DPM++          & \textbf{23.16}/40.23 & \textbf{9.74}/13.46 & \textbf{6.07}/6.83 & \textbf{4.81}/5.04 & \textbf{4.49}/\textbf{4.57} & \textbf{4.43}/\textbf{4.46}\\
DEIS           & \textbf{28.59}/35.24 & \textbf{11.62}/12.92 & \textbf{7.51}/\textbf{7.42} & \textbf{5.50}/\textbf{5.50} & \textbf{4.69}/\textbf{4.71} & \textbf{4.48}/\textbf{4.49}\\
UniPC          & \textbf{21.44}/40.41 & \textbf{9.18}/12.50 & \textbf{5.58}/6.06 & \textbf{4.66}/4.79 & \textbf{4.46}/\textbf{4.50} & \textbf{4.42}/\textbf{4.44}\\
\midrule
\multicolumn{7}{l}{\emph{Path family $P$ --- calibrated $m^\star$ vs.\ that path's uniform-grid default}}\\
1-RF   & \textbf{37.74}/\textbf{37.94} & \textbf{19.22}/19.74 & \textbf{12.28}/13.39 & \textbf{8.87}/10.26 & \textbf{6.75}/7.80 & \textbf{5.79}/6.35\\
2-RF   & \textbf{7.31}/7.54 & \textbf{6.80}/6.98 & \textbf{6.57}/6.69 & \textbf{6.44}/\textbf{6.51} & \textbf{6.33}/\textbf{6.37} & \textbf{6.27}/\textbf{6.29}\\
3-RF   & \textbf{7.30}/7.40 & \textbf{7.06}/7.16 & \textbf{6.95}/\textbf{7.02} & \textbf{6.88}/\textbf{6.93} & \textbf{6.83}/\textbf{6.85} & \textbf{6.79}/\textbf{6.80}\\
VP-SDE & \textbf{236}/240 & \textbf{131}/151 & \textbf{67.29}/80.13 & \textbf{37.06}/43.06 & \textbf{19.55}/22.28 & \textbf{12.32}/14.11\\
OT-CFM & 28.87/\textbf{27.86} & \textbf{17.32}/17.90 & \textbf{12.12}/13.52 & \textbf{9.39}/10.80 & \textbf{7.55}/8.47 & \textbf{6.67}/7.04\\
\midrule
\multicolumn{7}{l}{\emph{EMS solver --- genuine per-core $m^\star(d_{\mathrm{EMS}})$ vs.\ its own jointly co-designed schedule}}\\
EMS$^\ast$ & 19.33/\textbf{17.07} & 6.64/\textbf{6.30} & \textbf{4.71}/\textbf{4.71} & \textbf{4.38}/4.46 & \textbf{4.34}/\textbf{4.35} & ---$^\ddagger$\\
\bottomrule
\end{tabular}

\smallskip
{\footnotesize $^\dagger$Heun needs odd NFE ($2\,\mathrm{steps}{-}1$); both
arms ran $\{5,9,13,19,33,65\}$, i.e.\ $+1$ NFE per column. \emph{Correction:}
an earlier version of this row was invalid --- an NFE-accounting bug in the
runner gave the $m^\star$ arm $2\,\mathrm{NFE}{-}1$ true evaluations while
labeling them NFE (the default arm was unaffected; run manifests record the
true count). The cells shown are re-measured at matched \emph{true} NFE
through the fixed pipeline; the win/tie verdicts survive the correction with
smaller margins. $^\ast$EMS $=$ DPM-Solver-v3; ``default'' is its own EMS
schedule, co-derived \emph{jointly} with the solver from shared model
statistics (\S\ref{sec:solver}); our arm is the same solver under genuine
per-core $m^\star(d_{\mathrm{EMS}})$ at $k\!=\!2$, which a validation sweep
confirms is the optimal perceptual exponent for $d_{\mathrm{EMS}}$ --- a
sharp optimum ($k\!\le\!1$ under-concentrates, $k\!\ge\!3$ over-crowds the
nodes and the EMS solve diverges). $^\ddagger$Genuine $d_{\mathrm{EMS}}$ is
so low-$\sigma$-peaked (max/median $\approx195$ vs Heun's $4.1$) that at NFE
64 the inverse-CDF placement crowds near-identical nodes and the EMS
coefficient solve goes singular; all schedules tie at the floor there
regardless.}
\end{table}

\noindent\textbf{Reading the table.} Wins concentrate where a schedule has
room to act --- low-to-mid NFE, where node placement decides whether the
high-curvature region is resolved. As NFE grows every schedule converges to
the same floor and the cells become ties: this is the certificate's own
prediction (the discretization term $\int d/m^{p}\!\to\!0$), not a weakness.
The pattern is uniform across four solver cores and, crucially, across four
path families --- VE/EDM, the linear-interpolant RF, the VP-SDE, and the
optimal-transport-coupled CFM --- so $m^\star$ is not an EDM artifact. VP-SDE
is the cleanest case: Euler never reaches its floor in this NFE range, so
$m^\star$ wins at \emph{every} budget ($12$--$16\%$). Our schedule trails
only at the lowest budgets: the EMS solver at NFE 5--8 (against a schedule
co-optimized with the solver itself, yet our genuine per-core schedule still
ties at NFE$\,\ge\,$12 and edges ahead at 18), and OT-CFM at NFE 5 (too few
nodes to place on its sharply-peaked defect). OT-CFM is also an independent
witness to the \textbf{path} axis: its OT coupling beats the
independent-coupling 1-RF at low NFE (uniform $27.9$ vs $37.9$ at NFE 5) but
carries a slightly higher floor ($7.0$ vs $6.4$ at NFE 64) --- the same
$Q_0$-vs-$d$ tradeoff the reflow ladder shows (\S\ref{sec:path}), reached by
a different straightening mechanism.

\medskip
\noindent\textbf{Genuine per-core calibration matters.} Each core's
$m^\star$ is built from \emph{that core's own} measured defect $d_s$. For
EMS this required measuring $d_{\mathrm{EMS}}$ directly on the EMS trajectory
--- its multistep-plus-corrector update does not fit the single-step
calibration the other four cores share. An earlier version borrowed Heun's
$d_{\mathrm{Heun}}$; that gives a $17$--$43\%$ different grid and loses at
every NFE (e.g.\ $25.77$ vs $17.07$ at NFE 5). The genuine per-core defect
closes most of that gap (NFE 5: $19.33$) and turns the mid-NFE losses into
ties and a win. \emph{Per-core fairness needs the per-core defect, not merely
per-core node placement} --- the same principle that motivated calibrating
each core to itself in the first place.

\medskip
\noindent\textbf{Anatomy of one column} (1-RF, the cleanest single axis): the
calibrated grid's gain over uniform peaks mid-budget and is parameter-free.
\begin{table}[h]\centering
\caption{Schedule axis on 1-RF, in detail. Calibrated vs uniform; $\Delta<0$ = win.}
\label{tab:sched}
\begin{tabular}{lcccccc}
\toprule
NFE & 5 & 8 & 12 & 18 & 32 & 64\\
\midrule
uniform               & \textbf{37.94} & 19.74 & 13.39 & 10.26 & 7.80 & 6.35\\
proposed ($p\!=\!1$)  & \textbf{37.74} & \textbf{19.22} & \textbf{12.28} & \textbf{8.87} & \textbf{6.75} & \textbf{5.79}\\
$\Delta$              & $-0.20$ & $-0.52$ & $-1.11$ & $-1.39$ & $-1.05$ & $-0.55$\\
\bottomrule
\end{tabular}
\end{table}

\medskip
\noindent\textbf{$m^\star$ is FID-optimal, not merely better than the
default.} A free-node search (Nelder--Mead over all $K\!-\!1$ interior nodes,
seeded from \emph{both} $m^\star$ and uniform, optimized on a held-out seed
and validated on seeds 0/1/2 at 10k) cannot beat $m^\star$ on 1-RF:

\begin{center}\small
\begin{tabular}{lccc}
\toprule
1-RF & uniform & $m^\star$ & FID-optimal (search)\\
\midrule
NFE 8  & $19.74$ & $\mathbf{19.07}$ & $19.02$\\
NFE 12 & $13.39$ & $\mathbf{12.23}$ & $12.20$\\
\bottomrule
\end{tabular}
\end{center}

\noindent The search improves on $m^\star$ by only $0.04$--$0.05$ FID --- well
inside the 3-seed $2\sigma$ ($0.30$/$0.41$) --- while $m^\star$ beats uniform
by $0.66$/$1.16$. Crucially, seeding the search from \emph{uniform} converges
no lower than $m^\star$ at either budget, so this is not the optimizer merely
resting where it started. On 1-RF the closed form is thus statistically
indistinguishable from the FID-optimum of the entire free-node schedule
space. (A brute-force DP optimum of a discretized pixel-truncation surrogate
$\sum d\,\Delta t^{2}$ \emph{under}performs at $21.6$/$14.6$; with the caveat
that this DP uses a crude left-node cost, it still indicates the \emph{smooth}
certificate form, not naive surrogate minimization, is what tracks FID.)
This optimality is \emph{not} universal: the same search run on EDM/Heun
(Table~\ref{tab:schedaxis}) finds schedules $0.7$/$1.3$ FID below $m^\star$
at NFE 9/13. A full re-validation of $(p,k)$ at true NFE (18 configurations
$\times$ 5 budgets, validation seed only) \emph{rules out} mis-tuning as the
cause: the incumbent $(2,2)$ remains the global validation optimum ---
per-budget best outright at NFE 5--9 and within $0.1$ FID of best elsewhere
--- yet no member of the family reaches the search's schedules ($16.24$ vs
$15.55$ at NFE 9 on the same validation seed). The headroom is a genuine
limitation of the closed-form $(d\,\sigma^{-k})^{1/(p+1)}$ family at low NFE
on this cell --- the same finite-step regime where the bound's
lever-blindness binds --- not a tuning artifact.

\medskip
\noindent\textbf{The schedule axis in isolation (Table~\ref{tab:schedaxis}).}
Table~\ref{tab:dominance} compares $m^\star$ against each configuration's
\emph{native} default. Table~\ref{tab:schedaxis} asks the harder question:
within a fixed (path, solver) group, how does $m^\star$ rank against
\emph{optimized} schedulers --- a per-NFE tuned Karras $\rho$, the
published GITS dynamic program (Chen et al., 2024; reimplemented with the
group's own solver as jump operator), the theory-surrogate DP, and the
free-node FID search? Every row within a group is the identical sampler;
only node placement differs.

\begin{table}[t]\centering\small
\caption{\textbf{Schedule axis isolated.} Clean-FID (10k, 3 seeds) of
\emph{schedulers only}, within groups that fix path and solver; bold $=$
within $2\sigma$ of the column best in its group. Tuned $\rho$ and the
free-node search consume FID feedback per budget (and the search is run only
at the marked budgets); $m^\star$, GITS, and the surrogate-DP are metric-free.
The two $m^\star$ runs on 1-RF (independent calibrations) give
$19.07$--$19.22$ at NFE 8; the table shows the dominance-table run.}
\label{tab:schedaxis}
\setlength{\tabcolsep}{5pt}
\begin{tabular}{lcccccc}
\toprule
\multicolumn{7}{l}{\emph{EDM path + Heun core (odd NFE)}}\\[1pt]
\quad NFE & 5 & 9 & 13 & 19 & 33 & 65\\
\quad Karras $\rho\!=\!7$ (native default) & $338.52$ & $45.07$ & $11.64$ & $5.56$ & $\mathbf{4.45}$ & $\mathbf{4.40}$\\
\quad Karras tuned $\rho$ (per-NFE, val.-swept) & $245.02$ & $43.95$ & $9.98$ & $\mathbf{5.19}$ & $\mathbf{4.41}$ & $\mathbf{4.37}$\\
\quad GITS (DP, Heun-jump cost) & $242.52$ & $57.24$ & $18.07$ & -- & -- & --\\
\quad free-node FID search & -- & $\mathbf{12.70}$ & $\mathbf{6.43}$ & -- & -- & --\\
\quad $m^\star$ (closed form, $p\!=\!2$, $k\!=\!2$) & $\mathbf{46.22}$ & $13.41$ & $7.70$ & $5.33$ & $4.52$ & $\mathbf{4.41}$\\
\midrule
\multicolumn{7}{l}{\emph{RF path (1-RF) + Euler core}}\\[1pt]
\quad NFE & 5 & 8 & 12 & 18 & 32 & 64\\
\quad uniform-$t$ (native default) & $37.94$ & $19.74$ & $13.39$ & $10.26$ & $7.80$ & $6.35$\\
\quad DP of truncation surrogate $\sum d\,\Delta t^{2}$ & -- & $21.58$ & $14.60$ & -- & -- & --\\
\quad GITS (DP, Euler-jump cost) & $\mathbf{34.49}$ & $\mathbf{18.78}$ & $\mathbf{12.41}$ & $\mathbf{9.02}$ & $\mathbf{6.83}$ & $5.99$\\
\quad free-node FID search & -- & $\mathbf{19.02}$ & $\mathbf{12.20}$ & -- & -- & --\\
\quad $m^\star$ (closed form, $p\!=\!1$) & $37.74$ & $19.22$ & $\mathbf{12.28}$ & $\mathbf{8.87}$ & $\mathbf{6.75}$ & $\mathbf{5.79}$\\
\midrule
\multicolumn{7}{l}{\emph{RF path (2-RF) + Euler core}}\\[1pt]
\quad NFE & 5 & 8 & 12 & 18 & 32 & 64\\
\quad uniform-$t$ (native default) & $7.54$ & $6.98$ & $6.69$ & $\mathbf{6.51}$ & $\mathbf{6.37}$ & $\mathbf{6.29}$\\
\quad GITS (DP, Euler-jump cost) & $\mathbf{7.09}$ & $\mathbf{6.71}$ & -- & -- & -- & --\\
\quad $m^\star$ (closed form, $p\!=\!1$) & $7.31$ & $6.80$ & $\mathbf{6.57}$ & $\mathbf{6.44}$ & $\mathbf{6.33}$ & $\mathbf{6.27}$\\
\bottomrule
\end{tabular}
\end{table}

\noindent Three findings. \emph{(i) Tuning the default's knob does not
rescue it.} The best per-NFE $\rho$ (swept on validation, tested once)
closes almost none of the low-NFE gap --- $245$ vs $m^\star$'s $46$ at NFE 5,
$44$ vs $13$ at NFE 9: the failure of the Karras schedule at low budgets is
the one-parameter \emph{family's} shape, not the $\rho\!=\!7$ setting. Tuned
$\rho$ does edge past $m^\star$ at NFE 19--33 (e.g.\ $5.19$ vs $5.33$) ---
the honest price of being parameter-free, paid where all schedules are
within $\sim\!0.2$ of each other.
\emph{(ii) GITS is budget-aware but core-fragile.} On near-straight Euler
paths at tiny $K$ it beats $m^\star$ ($-2.7$ on 1-RF, $-0.2$ on 2-RF at NFE
5): its DP evaluates candidate jumps \emph{at the actual step lengths}, so it
sees what the certificate's bound cannot --- curvature near the \emph{end} of
a jump barely displaces the endpoint (lever $\int a(s)(t_{\mathrm{end}}-s)ds$;
$38\%$ of 1-RF curvature mass sits in $t\in[0.95,1]$ where the lever is
$\le\!0.05$). The measured endpoint deviations rank uniform\,/\,$m^\star$\,/\,GITS
exactly as FID does, and this lever-blindness is the \emph{same} mechanism
behind every low-NFE loss in Table~\ref{tab:dominance} (EMS at 5--8, OT-CFM
at 5). But the advantage does not survive a second-order core: with the
Heun-consistent jump cost, GITS collapses on EDM ($242/57/18$ vs $m^\star$'s
$46/13/8$) --- the teacher-forced jump metric badly misjudges $\sigma$-space
placement when the core itself cancels the leading curvature term. $m^\star$
is the only scheduler in the table that is never catastrophic in any group.
\emph{(iii) The remaining headroom is quantified.} The free-node search
bounds what any scheduler could achieve: $m^\star$ sits on that bound on
1-RF and $0.7$--$1.3$ FID above it on EDM/Heun at NFE 9--13.

\smallskip
\noindent\emph{Cost of the schedule.} ``Optimized'' baselines pay for their
gains; in net evaluations (the budget's own unit):

\begin{center}\footnotesize
\begin{tabular}{lccc}
\toprule
scheduler & preparation (net evals) & amortization & requires\\
\midrule
$m^\star$            & $1.8{\times}10^{4}$ & once per net & net only\\
GITS (Euler cost)    & $3.3{\times}10^{4}$ & once per net & net only\\
GITS (Heun cost)     & $2.4{\times}10^{6}$ & once per net & net only\\
EMS (DPM-v3)         & $\sim\!10^{6}$ & once per net & net $+$ training data\\
tuned $\rho$         & $3.6{\times}10^{5}$ \emph{per budget} & per (net, NFE) & net $+$ FID oracle\\
free-node search     & $\sim\!9{\times}10^{6}$ \emph{per budget} & per (net, NFE) & net $+$ FID oracle\\
\bottomrule
\end{tabular}
\end{center}

\noindent $m^\star$ is the cheapest member of the cheapest access class
(metric-free, data-free, closed-form per budget); the rows of
Table~\ref{tab:schedaxis} measure what each additional order of magnitude of
preparation --- or each additional information source --- actually buys.

\emph{Relation to learned schedules (AYS).} Align-Your-Steps
(Sabour et al., 2024) also optimizes schedules (via a KLUB objective) and reports gains
on CIFAR-10/EDM; its optimized CIFAR schedules are not published (only the
SD/SDXL/DeepFloyd/SVD schedules are released) and its FID protocol differs
(50k legacy vs our 10k Clean-FID), so a number-to-number comparison would
require reimplementing their optimizer. The free-node search above is the
stronger comparison on our protocol: it bounds what \emph{any} schedule
optimizer --- learned or otherwise --- could achieve at these budgets, and
$m^\star$ matches it without optimization, FID feedback, or training.

\emph{Adaptive step-size control is not a substitute.} A Dormand--Prince
RK45 solver (local-error-adaptive steps, the classical ``automatic
schedule'') on the same 1-RF ODE: it cannot operate below NFE ${\approx}19$
at all (each step costs 6 evaluations); at NFE 25 it is \emph{FID-blind} ---
FID is non-monotonic in the tolerance at matched NFE ($12.0$ at
tol $0.5$ vs $32.2$ at tol $0.1$), because the local error norm accepts a
single giant step across the high-curvature data end that is perceptually
catastrophic. It crosses the Euler frontier only at NFE ${\gtrsim}40$
($5.88$ at NFE 43, $5.29$ at NFE 83), where its \emph{fifth-order core} ---
an $s$-axis advantage, not a scheduling one --- dominates, mirroring EMS on
EDM. At deployment budgets, error-adaptive placement is exactly the
local/linear regime mismatch that $m^\star$'s calibrated global placement
avoids.

\medskip
\noindent\textbf{Cross-dataset transfer (FFHQ-64).} The recipe transfers to a
new dataset with \emph{nothing} re-tuned: on EDM's FFHQ-64 model we calibrate
$d$ on the FFHQ net (held-out seed, as always) and reuse the CIFAR-validated
exponents $(p,k)=(2,2)$ unchanged. Both arms are the identical Heun loop;
only the grid differs (FID: EDM's reference implementation against NVIDIA's
50k FFHQ statistics, same metric both arms; odd NFE; 10k, 3 seeds).
$m^\star$ beats the Karras $\rho\!=\!7$ default at \emph{every} budget,
including the highest ($\Delta=-0.10$ at NFE 33 exceeds $2\sigma=0.05$):

\begin{center}\small
\begin{tabular}{lccccc}
\toprule
FFHQ-64 (Heun) & NFE 5 & 9 & 13 & 19 & 33\\
\midrule
Karras $\rho\!=\!7$ default & $345.81$ & $58.50$ & $17.10$ & $6.40$ & $3.89$\\
$m^\star$ (CIFAR $(p,k)$, no re-tuning) & $\mathbf{37.17}$ & $\mathbf{11.33}$ & $\mathbf{6.34}$ & $\mathbf{4.39}$ & $\mathbf{3.79}$\\
\bottomrule
\end{tabular}
\end{center}

\noindent The FFHQ defect profile is genuinely different (max/median $2.2$
vs CIFAR's $4.1$ --- much flatter), so this is the calibration doing the
work, not a memorized CIFAR shape: the same closed form, fed a new $d$,
places nodes correctly on a new dataset.

\medskip
\noindent Two secondary findings (details in \S\ref{app:sched}):
\begin{itemize}\itemsep2pt
\item\emph{Path\,$\times$\,schedule:} the calibrated schedule wins on 2-RF
and 3-RF too, but the gain shrinks as reflow pushes the model toward its
floor (less defect to recover) --- straightening and scheduling are
complementary, not redundant.
\item\emph{FID-faithful weighting fails:} replacing $d$ with the
feature-weighted $d\!\cdot\!g$ is \emph{worse} than both pixel-$d$ and
uniform --- the local sensitivity $g$ is regime-mismatched. Pixel-$d$ is
sufficient and superior. (The EDM ``decomposability wall,'' cross-family.)
\end{itemize}

\section{Conclusion}
The certificate's structure holds on every axis we can probe. The two claims
that matter transfer off EDM: the \textbf{path} trades $Q_0$ for $d$ with an
NFE-dependent optimum (a confirmed crossover cascade), and the
\textbf{parameter-free calibrated schedule} wins or ties the native default
of every solver core and every path family tested --- four EDM cores and
four path families (EDM, RF, VP-SDE, OT-CFM) --- the only losses being at the
lowest budgets (EMS at NFE 5--8, where it still ties at NFE$\,\ge\,$12 with a
genuine per-core defect; OT-CFM at NFE 5; Table~\ref{tab:dominance}), and it
carries to a new dataset (FFHQ-64) with no re-tuning. Against
\emph{optimized} schedulers (Table~\ref{tab:schedaxis}) it is FID-optimal on
1-RF, within $\sim\!1$ FID of optimal on EDM/Heun, and uniquely
never-catastrophic --- at orders of magnitude less preparation compute and
the weakest access requirements (net only, no data, no metric).
Path-straightening,
schedule-calibration, and solver-core quality are complementary advances on
orthogonal axes of one certificate.

\appendix
\section{Protocol}
\label{app:protocol}
RF: NCSN++/DDPM++ backbone, EMA weights, standard Gaussian prior, Euler
$v$-ODE $x_{i+1}=x_i+v_\theta(x_i,t_i)(t_{i+1}-t_i)$, $t:10^{-3}\!\to\!1$,
NFE${}=K$. Clean-FID 10k, 3 seeds, $\pm$s.e.m. Floor probe (NFE 128/256)
uses the same integrator (no RK45 confound). Schedule calibration: $d(t)$
estimated once on a 512-step reference at a \emph{held-out} seed (777,
disjoint from eval seeds), nodes by equipartitioning $\int m^\star dt$.
Robustness: two independent defect estimators (second-difference
$\norm{x_{i+1}-2x_i+x_{i-1}}$ on the held-out seed; velocity-FD
$\norm{\dot v}$ on seed 0) agree (1-RF gap peaks $-1.39$ vs $-1.36$ at
$K\!=\!18$); the held-out-seed estimator is reported as primary.
FFHQ-64 transfer: EDM's pretrained \texttt{ffhq-64x64} model, $d_{\mathrm{Heun}}$
calibrated on the FFHQ net at the held-out seed, Heun loop identical in both
arms, odd NFE ($=2\,\mathrm{steps}-1$), 10k samples $\times$ 3 seeds, FID via
EDM's \texttt{fid.py} against NVIDIA's 50k FFHQ reference statistics.
Schedule-axis baselines (Table~\ref{tab:schedaxis}): tuned $\rho$ swept over
$\{2,3,4,5,7,10,15,20\}$ on a validation seed at 5k, best-per-NFE evaluated
once on test; GITS reimplemented with a 256-trajectory reference ensemble at
the held-out seed on a fine grid (128-node $t$ / 96-node $\log\sigma$), jump
cost from the group's own solver, DP per budget, tested once; free-node
search as on 1-RF (Nelder--Mead, search seed only, tested once).
\emph{NFE-accounting correction (2026-06-11):} the original Heun-row
$m^\star$ runs executed $2\,\mathrm{NFE}{-}1$ true evaluations under an NFE
label (runner bug, default arm unaffected); all run manifests were audited
(true count is recorded per run), the damage was confined to that row, and
the row was re-measured at matched true NFE through the fixed pipeline.
Cells whose original \emph{true} NFE already lay on the grid (9, 65) were
reused rather than re-drawn.

\section{Schedule supplement}
\label{app:sched}
\textbf{Exponent sensitivity} (seed-0 calibration, internally consistent),
proposed FID by exponent (bold = best per NFE):

\begin{center}\small
\begin{tabular}{lccccc}
\toprule
$K$ & 5 & 8 & 12 & 18 & 32\\
\midrule
$p\!=\!0.5$ & 42.71 & 20.99 & 12.97 & 9.16 & 6.86\\
$p\!=\!1$   & 37.17 & 19.07 & \textbf{12.23} & \textbf{8.90} & \textbf{6.79}\\
$p\!=\!2$   & \textbf{35.12} & \textbf{18.56} & 12.35 & 9.08 & 6.93\\
\bottomrule
\end{tabular}
\end{center}

\noindent The optimal exponent is mildly NFE-dependent ($p\!=\!2$ best at
$K\!\le\!8$, $p\!=\!1$ best at $K\!\ge\!12$, $p\!=\!0.5$ over-concentrates),
echoing the EDM finding that the deployment exponent is in-regime, not the
asymptotic order; $p\!=\!1$ is the theory-given parameter-free choice that
wins at \emph{every} NFE.

\medskip
\noindent\textbf{Path\,$\times$\,schedule gap grid} ($\Delta$FID = proposed
$-$ uniform, held-out calibration):

\begin{center}
\begin{tabular}{lcccccc}
\toprule
$\Delta$FID & $K{=}5$ & $K{=}8$ & $K{=}12$ & $K{=}18$ & $K{=}32$ & $K{=}64$\\
\midrule
1-RF & $-0.20$ & $-0.52$ & $-1.11$ & $-1.39$ & $-1.05$ & $-0.55$\\
2-RF & $-0.23$ & $-0.18$ & $-0.11$ & $-0.07$ & $-0.03$ & $-0.01$\\
3-RF & $-0.11$ & $-0.10$ & $-0.07$ & $-0.05$ & $-0.02$ & $-0.01$\\
\bottomrule
\end{tabular}
\end{center}

\noindent The defect \emph{shape} is constant along the ladder
($\norm{\ddot x}$ max/median $\approx69,72,72$ for 1/2/3-RF); reflow shrinks
its \emph{magnitude}, not its shape, so straighter paths run near their
floors and leave little FID room for scheduling --- consistent with
$m^\star$'s scale-invariance in $d$. (The $K\!=\!5$ 1-RF gap is anomalously
small: 5 nodes cannot tame its extreme curvature regardless of placement.)

\medskip
\noindent\textbf{FID-faithful weighting} (1-RF, seed-0). $g(t)=$ feature
sensitivity of the final sample to a perturbation at $t$; positive and
smooth but peaks at the noise end ($t\!\to\!0$). The $d\!\cdot\!g$ schedule
gives FID $48.4/26.3/16.9/11.6/8.0$ at $K\!=\!5/8/12/18/32$ --- worse than
uniform ($37.9/19.7/13.4/10.3/7.8$) and pixel-$d$
($37.2/19.1/12.2/8.9/6.8$): the local/linear $g$ over-weights the noise end
and starves the high-curvature data end.

\end{document}
